\section{Zadanie 9}
Podaj przykład zmiennej losowej, która ma skończoną wartość oczekiwaną,
i nieograniczoną wariancję.

\subsection*{Rozwiązanie}
W tym zadaniu chcemy znaleźć zbieżny szereg, który podniesiony do kwadratu jest rozbieżny. \\
Rozważmy następujący eksperyment: Losujemy liczbę ze zbioru liczb naturalnych.
Prawdopodobieństwo wylosowania liczby \(k\) jest równe \(\frac{C}{k^3}\).
\[ P(X=k) = \frac{C}{k^3} \]
\(C\) to tak dobrana stała żeby suma prawdopodobieństw wynosiła \(1\),
dokładniej \(C \approx 0,83\), ale to nas nie interesuje.

Zdefiniujmy zmienną losową \(X\) przyjmującą wartość wylosowanej liczby.

Z definicji wartości oczekiwanej: \\
\(E[X] = (1 \cdot P(X=1)) + (2 \cdot P(X=2)) + (3 \cdot P(X=3)) + \cdots\) \\
\(E[X] = (1 \cdot \frac{C}{1^3}) + (2 \cdot \frac{C}{2^3}) + (3 \cdot \frac{C}{3^3}) + \cdots\) \\
\(E[X] = C \cdot (\frac{1}{1^2} + \frac{1}{2^2} + \frac{1}{3^2} + \frac{1}{4^2} + \cdots)\)\\
a to z wielce znanego faktu jest skończone. OK.

To teraz wariancja:
z definicji \(Var(X) = E[X^2] - E[X]^2\). policzmy na początku \(E[x^2]\)\\
\(E[X^2] = (1^2 \cdot P(X=1)) + (2^2 \cdot P(X=2)) + (3^2 \cdot P(X=3)) + \cdots\) \\
\(E[X^2] = (1^2 \cdot \frac{C}{1^3}) + (2^2 \cdot \frac{C}{2^3}) + (3^2 \cdot \frac{C}{3^3}) + \cdots\) \\
\(E[X^2] = C \cdot (\frac{1}{1} + \frac{1}{2} + \frac{1}{3} + \frac{1}{4} + \cdots)\)\\
co ciekawe znowu otrzymaliśmy znany szereg, tym razem nieskończony.
Skoro \(E[X^2]\) jest nieskończone, a \(E[X]\) skończone, to \(E[X^2] - E[X]^2\) (czyli wariancja) jest nieskończona.

\subsection*{Odpowiedź}
\[
	\boxed{\text{Jest taki przykład, jak najbardziej, pokazany w rozwiązaniu}}
\]
